%%=============================================================================
%%===   ISEA-TEX-Vorlage   ===
%%=== Formblatt zur Arbeit ===
%%===        mla 2016      ===
%%=============================================================================

%%	Diese Datei braucht i.d.R. nicht editiert zu werden, es werden die Werte aus isea_tex_def übernommen.

\label{SummaryForm}

\thispagestyle{empty}
%\edef\pagenumstart{\value{page}}

\small
Diese Seite wird nicht mitgebunden und ist als separate Seite bei der Abgabe der Arbeit mit abzugeben.
%\normalsize
\vspace{5mm}

\varTitle
\vspace{5mm}

\varDate
\vspace{5mm}

Autor: \varAuthorFamily{}, \varAuthorGiven{}

Betreuende/r Mitarbeiter/in: \varAssiName

Betreuender Professor: \varSupA

\vspace{5mm}
\ifthenelse{\equal{\varInstitute}{ISEA}}
{
	Arbeit außerhalb des ISEA angefertigt: \textbf{\ifthenelse{\equal{\varExtern}{true}}{ja}{nein}}
}
{
	Arbeit außerhalb des PGS angefertigt: \textbf{\ifthenelse{\equal{\varExtern}{true}}{ja}{nein}}
}

\ifthenelse{\equal{\varExtern}{true}}
{
	Firma der Anfertigung: \varExternName
}
{}

\vspace{5mm}

Studienrichtung: \varFieldOfStudy

Sprache: \ifthenelse{\equal{\varLanguage}{german}}{Deutsch}{Englisch}

\vspace{7mm}
Vertraulichkeit:

\ifthenelse{\equal{\varInstitute}{ISEA}}
{
	$\square$ Diese Arbeit fällt unter eine Geheimhaltungsvereinbarung des ISEAs und beinhaltet vertrauliche Daten.
}
{
	$\square$ Diese Arbeit fällt unter eine Geheimhaltungsvereinbarung des PGS und beinhaltet vertrauliche Daten.	
}

$\square$ Diese Arbeit steht unter erweiterten Geheimhaltungsvereinbarungen mit Dritten oder zum Schutz von Erfindungsrechten.

$\square$ Diese Arbeit hat einen Sperrvermerk bis zum \ldots \ldots \ldots .

\vspace{5mm}

Stichworte:
\begin{multicols}{2}
	\begin{itemize}
		\itemsep 0pt
		\makeatletter
		\@for\tmp:=\varKeywords\do{\item \tmp}
		\makeatother
	\end{itemize}
\end{multicols}


\vspace{5mm}

Zusammenfassung:


Hier steht eine kurze Zusammenfassung der Arbeit. Der Umfang der Zusammenfassung sollte vier bis acht Zeilen betragen.

Dieser Text ist in der Datei \textbf{SummaryForm.tex} zu finden, weitere Parameter (Betreuer, Studienrichtung, ...) können entweder in der \textbf{isea\_tex\_def.tex} oder direkt in dieser Datei angepasst werden.


\cleardoubleemptypage
\normalsize
%\setcounter{page}{\pagenumstart-1}	% Zähler zurücksetzen, als ob Seite nicht vorhanden wäre
